\begin{textsample}{POS Dim 3 – human – Score 7.00 – rj/rj\_ef\_ciencias\_humanas\_his\_1\_p4.txt}  \label{ex:f3_pos_003}
A partir da apreciação de \textbf{especificidades} apontadas em \textbf{diretrizes} \textbf{curriculares} de vários municípios do Rio de Janeiro que vêm se debruçando sobre as premissas \textbf{colocadas} na BNCC, buscamos registrar aquelas que vão dando ao \textbf{Documento} \textbf{Curricular} do \textbf{Estado} do Rio de Janeiro uma característica própria, que vai além do proposto a nível \textbf{nacional}, sem, no entanto, descaracterizar um \textbf{documento} que já é um marco para a educação \textbf{nacional}.
Muito ao contrário, ao buscar \textbf{descrever} contextos próprios em habilidades específicas, a \textbf{proposta} do \textbf{Documento} \textbf{Curricular} do \textbf{Estado} do Rio de Janeiro vai \textbf{trazendo} mais concretude para as “ pistas ” apresentadas na BNCC e que se desdobrarão em estratégias próprias de cada Projeto Político-Pedagógico, em cada escola e sala de aula.

% matched lemmas: colocar, curricular, descrever, diretriz, documento, especificidade, estado, nacional, proposta, trazer
\end{textsample}
